\documentclass[conference]{IEEEtran}
\usepackage[T1]{fontenc}
\usepackage{amsmath}
\usepackage{graphicx}

\title{Capacitor Charging and a Power On Reset
for an Active Low Microcontroller}

\author{
\IEEEauthorblockN{Pina Cammisano}
\IEEEauthorblockA{Department of Computer Engineering Technology\\
Vanier College}
}

\begin{document}
\maketitle

\begin{abstract}
This document introduces the charging behaviour of a resistor-capacitor
(RC) network and applies the same charging behaviour to a power-on-reset
circuit for an active-low microcontroller.
\end{abstract}


\section{Introduction}

When a capacitor $C$ is charged through a series resistor $R$ from a
constant supply $V_0$, the capacitor voltage rises according to

\begin{equation}
V(t)=V_0\left(1-e^{-t/RC}\right)
\label{eq:charging}
\end{equation}

The product $RC$ is called the time constant,

\begin{equation}
\tau=RC.
\label{eq:tau}
\end{equation}

After one time constant, the capacitor reaches approximately 63.2\%
of the supply voltage.

\begin{figure}[t]
\centering
\includegraphics[width=\columnwidth]
{figures/generated/rc_charging.pdf}
\caption{Capacitor charging for different resistance values.}
\label{fig:rc}
\end{figure}


\section{Voltage at the Time Constant}

Starting from \eqref{eq:charging} and substituting $t=\tau$,

\begin{align}
V(\tau)
&=V_0\left(1-e^{-\tau/RC}\right) \\
&=V_0\left(1-e^{-RC/RC}\right) \\
&=V_0\left(1-e^{-1}\right) \\
&\approx V_0(1-0.368) \\
&\approx0.632V_0.
\end{align}

For $V_0=5.0$ V,

\begin{equation}
V(\tau)\approx0.632(5.0)\approx3.16\text{ V}.
\label{eq:vtau}
\end{equation}


\section{Power On Reset with an Active Low Reset}

A power-on-reset circuit holds the microcontroller in reset while the
capacitor charges. The active-low reset pin remains low while the
capacitor voltage is below the input-high threshold $V_{IH}$. Once the
voltage crosses $V_{IH}$, the reset is released and the microcontroller
can begin operating.

\begin{figure}[t]
\centering
\includegraphics[width=\columnwidth]
{figures/drawn/por_circuit.pdf}
\caption{RC power-on-reset circuit for an active-low reset pin.}
\label{fig:por}
\end{figure}

The reset pin releases when the capacitor voltage reaches $V_{IH}$:

\begin{equation}
V_{IH}=V_0\left(1-e^{-t_{\mathrm{release}}/RC}\right).
\end{equation}

Dividing by $V_0$,

\begin{equation}
\frac{V_{IH}}{V_0}
=
1-e^{-t_{\mathrm{release}}/RC}.
\end{equation}

Rearranging,

\begin{equation}
e^{-t_{\mathrm{release}}/RC}
=
1-\frac{V_{IH}}{V_0}.
\end{equation}

Taking the natural logarithm,

\begin{equation}
-\frac{t_{\mathrm{release}}}{RC}
=
\ln\left(1-\frac{V_{IH}}{V_0}\right).
\end{equation}

Therefore,

\begin{equation}
t_{\mathrm{release}}
=
-RC\ln\left(1-\frac{V_{IH}}{V_0}\right).
\label{eq:release}
\end{equation}

Solving for $R$ gives

\begin{equation}
R=
\frac{-t_{\mathrm{release}}}
{C\ln\left(1-\frac{V_{IH}}{V_0}\right)}.
\label{eq:rvalue}
\end{equation}

For $V_{IH}=0.7V_0$, $t_{\mathrm{release}}=10$ ms, and
$C=100$ nF,

\begin{equation}
\ln(1-0.7)=\ln(0.3)\approx-1.204.
\end{equation}

Using \eqref{eq:rvalue},

\begin{equation}
R=
\frac{-(10\times10^{-3})}
{(100\times10^{-9})(-1.204)}
\approx83\,\text{k}\Omega.
\end{equation}

An $83\,\text{k}\Omega$ resistor with a $100$ nF capacitor gives a
reset hold-off of approximately 10 ms.


\section{Applications and Limitations}

\subsection{Application}

A low-cost temperature sensor node is a suitable application for an
RC-only power-on-reset circuit. It only needs a short delay when power
is first applied and does not require precise reset timing, brown-out
detection, or re-triggering.

For this application,

\begin{equation}
R=10\,\text{k}\Omega,\qquad C=220\,\text{nF}.
\end{equation}

Assuming $V_{IH}=0.7V_0$ and using \eqref{eq:release},

\begin{equation}
t_{\mathrm{release}}
=-(10\times10^3)(220\times10^{-9})\ln(0.3)
\approx2.65\,\text{ms}.
\end{equation}


\begin{figure}[!t]
\centering
\includegraphics[width=\columnwidth]
{figures/generated/rc_tolerance.pdf}
\caption{Nominal, fastest, and slowest RC charging curves.}
\label{fig:tolerance}
\end{figure}


\subsection{Limitations}

One limitation is component tolerance. The actual resistor and
capacitor values can vary from their nominal values, which changes the
reset timing.

A second limitation is a slowly rising supply voltage. If $V_{CC}$
rises slowly, the reset pin can remain near its threshold and may not
produce a clean reset.

A third limitation is that the RC circuit does not provide brown-out
detection. If the supply voltage drops after startup, the circuit may
not reset the microcontroller.


\subsection{Component Tolerance}

For this design, a $10\,\text{k}\Omega$ resistor with a tolerance of
$\pm1\%$ and a $220\,\text{nF}$ capacitor with a tolerance of
$\pm10\%$ are used.

Since \eqref{eq:release} can be written as

\begin{equation}
t_{\mathrm{release}}=kRC,
\end{equation}

where

\begin{equation}
k=-\ln\left(1-\frac{V_{IH}}{V_0}\right),
\end{equation}

the effect of changes in $R$ and $C$ can be written as

\begin{align}
t+\Delta t
&=k(R+\Delta R)(C+\Delta C) \\
&=kRC+kR\Delta C+kC\Delta R+k\Delta R\Delta C.
\end{align}

Ignoring the second-order term,

\begin{equation}
t+\Delta t
\approx
kRC+kR\Delta C+kC\Delta R.
\end{equation}

Dividing by $t=kRC$ gives

\begin{equation}
\frac{\Delta t}{t}
\approx
\frac{\Delta R}{R}
+
\frac{\Delta C}{C}.
\label{eq:tolerance}
\end{equation}

Using the component tolerances,

\begin{equation}
\frac{\Delta t}{t}
\approx0.01+0.10=0.11.
\end{equation}

Therefore, the approximate timing tolerance is $\pm11\%$.

Using the nominal release time of $2.65$ ms,

\begin{equation}
t_{\min}\approx2.65(1-0.11)\approx2.36\,\text{ms},
\end{equation}

and

\begin{equation}
t_{\max}\approx2.65(1+0.11)\approx2.94\,\text{ms}.
\end{equation}

Therefore, the reset release time is approximately between
$2.36$ ms and $2.94$ ms.

\end{document}
